\documentclass[12pt]{article}
\usepackage[utf8]{inputenc}
\usepackage[T1]{fontenc}
\usepackage{hyperref}
\title{Ternate\~no Language Resources for K-12 Curriculum Development}
\author{Compiled by Claude Sonnet 5 (claude-sonnet-5) for Andr\'as Kornai}
\date{}
\begin{document}
\maketitle

\section{Introduction}
Ternate\~no (ISO 639-3 tmg, Glottocode tern1248), also called Portugis, is an extinct Portuguese-based creole. Christians of mixed Portuguese and Malay descent spoke it on Ternate and Ambon in the Moluccas, and Ambonese Malay gradually replaced it. Do not confuse it with Ternate Chabacano (Chavacano), a living Spanish-based creole of Ternate, Cavite, Philippines. Its origin is traced to the Moluccan island of Ternate. The two are related historically, and Verbix suggests Ternate\~no may have contributed to Ternate Chabacano. The searches found no grammars, dictionaries or recordings of Ternate\~no itself.

\section{Primary Educational Resources}
None were found.

\section{Linguistic and Reference Materials}
\begin{itemize}
\item Portugis (Wikipedia): \url{https://en.wikipedia.org/wiki/Portugis}
\item Verbix page for Ternate\~no: \url{https://docs.verbix.com/Languages/PortugueseTernatenyo}
\item Verbix page for Ternate Chabacano (the related Philippine creole): \url{https://docs.verbix.com/Languages/ChavacanoTernate}
\item APICS survey chapter on Ternate Chabacano (the Philippine creole, not Ternate\~no): \url{https://apics-online.info/surveys/44}
\end{itemize}

\section{Religious and Community Resources}
None were found.

\section{Digital and Audio Resources}
None were found.

\section{Curriculum Development Context}
Ternate\~no has no speakers, and the searches found no teaching materials. A heritage curriculum could cover the history of the Moluccan Portuguese-creole community and its link to Ternate Chabacano in the Philippines, and could draw on the living Chabacano communities. A language-revitalization component would be speculative without primary documentation.
\end{document}